\begin{table}[t]
\tbl{\label{tab:orbits:1}}
{\begin{tcolorbox}[tabularx={p{3cm}|p{2cm}|p{6.3cm}|},enhanced,fonttitle=\bfseries\large,fontupper=\normalsize\sffamily,
colback=yellow!10!white,colframe=red!50!black,colbacktitle=blue!30!white,
coltitle=black,title=Computing Orbits]
\hline
Command & Arguments & Purpose \\
\hline
\hline
\verb'-on_points'  &     &  Perform Schreier algorithm to compute the orbits of $G$ in the natural action and under the group included in the action. \\
\hline
\verb'-on_points_with_generators'  &  gens   &  Perform Schreier algorithm to compute the orbits of $G$ in the natural action and under the group given by a set of generators. \\
\hline
\verb'-on_subsets'  &  $k$ control  &  Consider the induced action on subsets of size $k$. This option will invoke the poset classification algorithm. \\
\hline
\verb'-of_one_subset'  &  set-label  &  Compute the orbit of the given set under the group $G$. This algorithm will build up the Schreier tree. The orbit will be saved to a csv file. \\
\hline
\verb'-on_subspaces'  &  $k$ control  &  Consider the induced action on subspaces of size $k$. We assume that the group action is linear. This option will invoke the poset classification algorithm. \\
\hline
\verb'-on_tensors'  &  $d$ control   &  Compute the orbits on the $d$-fold tensor product. \\
\hline
\verb'-on_partition'  &   $k$ control  &  Compute the orbits on partitions of type $k+k+k$. In this case, $3k$ must equal the degree of the defining action of $G$. This option will invoke the poset classification algorithm. \\
\hline
\verb'-on_polynomials'  &   R  &  Consider the action of $G$ on homogenous polynomials as defined in ring R. The action must be linear and the number of variables in R must match the degree of the matrix group. This option will invoke the Schreier orbit algorithm. \\
\hline
\verb'-on_one_polynomial'  &   R equation  &  Consider the action of $G$ on polynomials of degree $d$ in $n$ variables. $G$ must be a matrix group. Here, the action must be linear on an $n$-dimensional space. \\
\hline
\verb'-on_cubic_curves'  &   options  &  Classify cubic curves by means of $(9,3)$-arcs. For valid options, see Tab.~\ref{tab:arcs}. \\
\hline
\verb'-on_cubic_surfaces'  &   P control  &  Classify cubic surfaces in the projective space P, using the given poset classification control options. \\
\hline
\verb'-on_arcs'  &   P control  &  Classify arcs in the projective space P, using the given arc classification control options, see Tab.~\ref{tab:arcs}. This option will invoke the poset classification algorithm. \\
\hline
\verb'-classification_by_canonical_form'  &   description  &  Perform the classification-by-canonical-form algorithm using Nauty. See Chapter~\ref{chapter:canonical:forms}. \\
\hline
\end{tcolorbox}}
\end{table}
